\documentclass[11pt]{article}
\usepackage[margin=1in]{geometry}
\usepackage{amsmath,amssymb,amsthm}
\usepackage{graphicx}
\usepackage{booktabs}
\usepackage{hyperref}
\usepackage{natbib}
\usepackage{xcolor}
\usepackage{microtype}
\usepackage{multirow}

\newcommand{\Gr}{\mathrm{Gr}}
\newcommand{\BN}{\mathrm{BN}}
\newcommand{\R}{\mathbb{R}}
\newcommand{\TODO}[1]{\textcolor{red}{[TODO: #1]}}
\newcommand{\CITE}[1]{\textcolor{red}{[CITE: #1]}}
\newcommand{\dg}{d_{\mathrm{g}}}
\newcommand{\effdim}{\mathrm{dim}_{\mathrm{eff}}}

\title{Composition-Controlled Evaluation of RNA Structure Awareness\\
  in Foundation Models}

\author{Elliot Tower\\
  Mechanistic Research\\
  \texttt{elliot@elliottower.ai}}

\date{July 2026}

\begin{document}
\maketitle

\begin{abstract}
We introduce a composition-controlled evaluation framework for
measuring RNA secondary structure awareness in foundation models and
apply it across four preregistered phases to ten models spanning
attention-based, SSM-based, and hybrid architectures across four
orders of magnitude in parameter count (1.2M--7B).
The framework addresses a composition confound inherent to RNA
secondary structure: stems are 68.5\% GC while loops are 43.9\%, so
any embedding metric sensitive to nucleotide identity inherits a
${\sim}3\times$ bias toward stem positions.  Our nucleotide-stratified
permutation null absorbs this confound, reducing the apparent
structure signal from $5.5\times$ to $1.8\times$ (mean across 52
families).  A max-over-layers correction ensures the null is computed
under the same best-layer selection as the real data.

The framework reveals a comprehensive null: no model encodes secondary
structure beyond what composition and architecture explain.  NT~v2
(ESM+GLU, 56M, DNA-pretrained) shows high attention-contact
correlation ($\rho = 0.322$), but untrained weights produce the same
signal ($\rho = 0.328$), exposing this as an architectural artifact of
the ESM attention mechanism.  Phase~3 adds two RNA-pretrained models:
RiNALMo (650M) shows genuine but modest mutation sensitivity (17/52
families exceed the composition null, sign test 94\%), while UTR-LM
(1.2M) shows no structure awareness despite structure-aware
pretraining objectives.  Phase~4 adds three more models: ERNIE-RNA
(86M, RNA-pretrained) and SpliceBERT (19.4M, RNA-pretrained) both
confirm the null pattern on mutation and attention; DNABERT-2 (117M,
DNA, BPE tokenization) shows the only clear learned attention signal
($\rho = 0.257$ trained vs $0.000$ untrained), though this excludes
ALiBi position biases.  Evo (StripedHyena, 7B) shows the weakest
mutation sensitivity despite being $500\times$ larger than Caduceus
(14M), indicating that pretraining domain and tokenization dominate
scale.  No family survives both nucleotide and dinucleotide
composition controls.  The directional effect is real---trained
exceeds untrained in 50/52 families for RNA-FM and 52/52 for
ERNIE-RNA---but the magnitude remains below the preregistered
threshold for practical significance.

The framework generalizes to any domain where feature composition
correlates with the target structure.
\end{abstract}


%% ────────────────────────────────────────────
\section{Introduction}
\label{sec:intro}

Foundation models pretrained on nucleotide sequences have been
proposed for RNA biology, but evaluating whether they encode secondary
structure---the pattern of intramolecular base pairing that governs
RNA stability and function---requires metrics that distinguish learned
representations from sequence-composition artifacts.

We initially adopted the Grassmannian geodesic distance ($\dg$), a
geometric metric for comparing embedding subspaces between wild-type
and mutant sequences.  A zero-parameter 3-mer frequency baseline
matched the trained models' $\dg$ scores, falsifying the metric: it
captured $k$-mer composition, a property that random linear maps
preserve under the Johnson--Lindenstrauss lemma
\citep{johnson1984extensions}, rather than learned structure.  Similar
failure modes have been reported for geometric transfer metrics in
other domains \citep{tower2026boundary}.

The $\dg$ collapse motivated a replacement metric with explicit
composition controls.  We designed a mutation sensitivity ratio under
complement swap (A$\leftrightarrow$U, C$\leftrightarrow$G) with a
nucleotide-stratified permutation null.  The key observation: stems
are enriched for G and C (68.5\% vs 43.9\% in loops), and complement
swaps at G/C positions produce inherently larger embedding
perturbations than at A/U positions regardless of structural context.
Any naive stem-vs-loop comparison inherits this 24.7 percentage-point
enrichment as a confound.

The contribution is the evaluation framework and an honest report of
what it reveals: a comprehensive null on structure encoding across
ten models, an informative scaling observation (7B parameters do not
compensate for domain mismatch), and a corrected attention analysis
showing that NT~v2's high attention-contact correlation is
architectural rather than learned.  Structure probing of RNA models has previously focused on
downstream task performance \citep{rnafm,dalla2023nucleotide} or
attention visualization rather than composition-controlled metrics.
\citet{ethayarajh2019} showed that BERT embeddings become increasingly
anisotropic through layers, a phenomenon we observe in RNA-FM on
repeat sequences.  Vig et al.\ demonstrated that protein language
model attention heads correlate with contact maps
\citep{vig2020bertology}; our attention-contact analysis extends this
to RNA, finding strong learned correlation for NT~v2 and near-zero
for RNA-FM.  The mutation sensitivity approach relates to \emph{in
silico} mutagenesis for variant effect prediction
\citep{frazer2021disease}, applied here to the representation layer
with structure-defined mutation categories (stem vs.\ loop).
Composition confounds in genomic analyses have a long history:
dinucleotide frequencies affect codon usage statistics
\citep{karlin1998dinucleotide}, CpG content confounds methylation
studies, and GC content correlates with gene density and chromatin
accessibility.  Our stratified permutation framework adapts these
insights to the evaluation of neural sequence models.


%% ────────────────────────────────────────────
\section{Methods}
\label{sec:methods}

\subsection{RNA families and models}
\label{sec:families_models}

Phase~1 evaluated 12 RNA families with experimentally determined
secondary structures from PDB crystal structures and Rfam consensus
annotations (Table~\ref{tab:rna_families}).  Phase~2 expands to
$N = 52$ families drawn from Rfam~14.10, spanning ribozymes,
riboswitches, snRNAs, cis-regulatory elements, miRNA precursors,
CRISPR repeats, and other ncRNAs (Appendix~A).  Positions are labeled
stem (base-paired) or loop (unpaired) from dot-bracket annotation.
Minimum criteria: $\geq 20$~nt, $\geq 5$ stem and 5 loop positions.

\begin{table}[t]
\centering
\caption{Phase~1 RNA families.  Source: PDB crystal structures and
  Rfam consensus annotations.  Stem: number of base-paired positions.}
\label{tab:rna_families}
\small
\begin{tabular}{@{}llccl@{}}
\toprule
\textbf{RNA family} & \textbf{Class} & \textbf{Length} &
  \textbf{Stem} & \textbf{Source} \\
\midrule
tRNA$^{\mathrm{Phe}}$ (yeast)  & tRNA         & 76  & 42 & PDB 1EHZ \\
tRNA$^{\mathrm{Ala}}$ (human)  & tRNA         & 76  & 42 & Rfam RF00005 \\
5S rRNA (\textit{E.\ coli})    & rRNA         & 120 & 72 & PDB 1C2X \\
Hammerhead ribozyme             & Ribozyme     & 48  & 16 & PDB 2OEU \\
SAM riboswitch                  & Riboswitch   & 94  & 52 & PDB 2GIS \\
TPP riboswitch                  & Riboswitch   & 80  & 40 & PDB 2CKY \\
SRP RNA helix~8                 & Structural   & 77  & 44 & PDB 1HQ1 \\
miR-21 precursor                & Regulatory   & 72  & 64 & Rfam RF00001 \\
IRES HCV domain~II              & Regulatory   & 77  & 50 & PDB 1P5O \\
HDV ribozyme                    & Ribozyme     & 76  & 48 & PDB 1DRZ \\
RNase~P specificity             & Structural   & 78  & 42 & PDB 1NBS \\
U2 snRNA stem                   & Spliceosomal & 76  & 46 & PDB 3K1V \\
\bottomrule
\end{tabular}
\end{table}

Phases~1--2 evaluate five models spanning attention-based, SSM-based,
and hybrid architectures (Table~\ref{tab:models}): RNA-FM (BERT, 99M,
character, RNA-pretrained; \citealp{rnafm}), NT~v2 (ESM+GLU, 56M,
6-mer, DNA; \citealp{dalla2023nucleotide}), HyenaDNA (Hyena SSM,
5.4M, character, DNA; \citealp{nguyen2024hyenadna}), Caduceus
(BiMamba, 14M, character, DNA; \citealp{schiff2024caduceus}), and Evo
(StripedHyena, 7B, byte-level, DNA; \citealp{nguyen2024evo}).
Phase~3 adds RiNALMo (BERT+RoPE+SwiGLU, 650M, character, RNA;
\citealp{penic2024rinalmo}) and UTR-LM (transformer encoder, 1.2M,
character, UTR-focused; \citealp{chu2024utrlm}).
Phase~4 adds ERNIE-RNA (BERT, 86M, character, RNA;
\citealp{wang2024ernierna}), SpliceBERT (BERT, 19.4M, character, RNA;
\citealp{chen2024splicebert}), and DNABERT-2 (MosaicBERT+ALiBi, 117M,
BPE, DNA; \citealp{zhou2024dnabert2}).  Six of the ten models are
RNA-pretrained; parameter counts span four orders of magnitude.

\begin{table}[t]
\centering
\caption{Models under test.  Phases~1--2 evaluate the first five;
  Phase~3 adds RiNALMo and UTR-LM; Phase~4 adds the final three.}
\label{tab:models}
\small
\begin{tabular}{@{}llcll@{}}
\toprule
\textbf{Model} & \textbf{Architecture} & \textbf{Params} &
  \textbf{Tokenization} & \textbf{Pretraining} \\
\midrule
RNA-FM     & BERT encoder          & 99M  & Character  & RNA \\
NT~v2      & ESM + GLU attention   & 56M  & 6-mer      & DNA \\
HyenaDNA   & Hyena SSM             & 5.4M & Character  & DNA \\
Caduceus   & BiMamba SSM           & 14M  & Character  & DNA \\
Evo        & StripedHyena          & 7B   & Byte-level & DNA \\
\midrule
RiNALMo    & BERT + RoPE + SwiGLU  & 650M & Character  & RNA \\
UTR-LM     & Transformer encoder   & 1.2M & Character  & UTR (RNA) \\
\midrule
ERNIE-RNA  & BERT encoder          & 86M  & Character  & RNA \\
SpliceBERT & BERT encoder          & 19.4M & Character & RNA \\
DNABERT-2  & MosaicBERT + ALiBi    & 117M & BPE        & DNA \\
\bottomrule
\end{tabular}
\end{table}

\subsection{Structure sensitivity metrics}
\label{sec:metrics}

\textbf{Mutation sensitivity ratio.}  For each position in each RNA
sequence, we perform a complement swap (A$\leftrightarrow$U,
C$\leftrightarrow$G) and measure the cosine distance between
wild-type and mutant embeddings at the mutated position.  The
structure sensitivity ratio $R$ is the mean cosine distance at stem
positions divided by the mean at loop positions.  We compute $R$ at
every layer and report the maximum; the null models below are computed
under the same max-over-layers selection.

\textbf{Attention-contact correlation.}  For models with accessible
attention matrices (RNA-FM, NT~v2, RiNALMo, UTR-LM, ERNIE-RNA,
SpliceBERT, DNABERT-2), we compute the Spearman rank correlation
between symmetrized attention weights and a binary contact matrix (1
for base-paired positions, 0 otherwise).  For NT~v2's 6-mer and
DNABERT-2's BPE tokenization, contacts are aggregated to the token
level.  All models except DNABERT-2 are loaded with
\texttt{attn\_implementation=``eager''} to ensure attention weights
are returned.  DNABERT-2's MosaicBERT architecture uses a fused Wqkv
projection; we extract content-only attention by hooking the Wqkv
output, splitting Q and K, and computing
$\mathrm{softmax}(QK^T/\sqrt{d_k})$ without ALiBi position biases.
We compare trained attention to an untrained baseline with
randomized weights.

\textbf{Structure probing.}  Logistic regression on per-position
embeddings predicting stem vs loop, using GroupKFold cross-validation
by RNA family to test cross-family generalization.  Balanced accuracy
is reported against the majority baseline.

\subsection{Composition null models}
\label{sec:nulls}

The nucleotide-stratified null shuffles stem/loop labels independently
within each nucleotide type (A, U, G, C), preserving per-nucleotide
composition while destroying the structure-label association.  For
each of 100 permutations, a single shuffled label assignment is
generated and $R$ is computed at every layer; the maximum across
layers is taken.  The 95th percentile of these 100 max-over-layers
values defines the null threshold.  This controls for first-order
composition (per-nucleotide GC enrichment) but not for dinucleotide
context---stem-G occurs predominantly in GC base pairs, while loop-G
occurs in diverse dinucleotide contexts.

The dinucleotide-stratified null extends this to second-order context:
labels are shuffled within each dinucleotide type (previous nucleotide
+ current nucleotide: AA, AC, AG, \ldots).  The same max-over-layers
procedure applies.  This null is run on all families exceeding the
first-order null, plus tRNA-Ala and SAM riboswitch (Phase~1 survivors)
regardless of Phase~2 result.

Both phases were preregistered with frozen code and analysis SHAs
(Phase~1: \texttt{694b43b}; Phase~2: \texttt{bd4b3fd}).  Phase~1
($N = 12$) tested whether RNA-FM's trained/untrained ratio exceeds
2.0 (H1), whether $\geq 6/12$ families exceed the null (H1b), and
whether probing exceeds the majority baseline (H5).  Phase~2
($N \geq 50$) replicated these at scale (H6, H6b, H8), added a
dinucleotide null test (H7: $\geq 1$ family survives), a confirmatory
sign test promoting the Phase~1 exploratory unanimity (H10: trained
$>$ untrained in $\geq 75\%$ of families), and an NT~v2 attention
test (H11: trained $\rho > 0.15$, untrained $< 0.05$).  A
methodological correction applied to both phases: the Phase~1 null
had been computed per-layer independently, while the metric used
max-over-layers.  Phase~2 corrects this mismatch.


%% ────────────────────────────────────────────
\section{Results}
\label{sec:results}

\subsection{Attention-contact correlation is architectural}
\label{sec:attn_results}

NT~v2 trained attention correlates with base-pairing contacts at mean
Spearman $\rho = 0.322$ across 52 families
(Table~\ref{tab:attention}).  An earlier analysis using the default
SDPA attention backend reported $\rho = 0.000$ for untrained weights,
but SDPA silently returns \texttt{None} for attention weights,
producing a spurious zero.  With eager attention computation,
untrained NT~v2 produces $\rho = 0.328$---slightly \emph{higher} than
trained.  The per-family delta (trained $-$ untrained) averages
$-0.006$ with no systematic direction: 22 families show higher trained
correlation, 30 show higher untrained.  The entire attention-contact
signal is an architectural property of the ESM attention mechanism and
6-mer tokenization, present from random initialization.

\begin{table}[t]
\centering
\caption{Attention-contact Spearman correlation ($N = 52$).  NT~v2's
  high correlation is entirely architectural (untrained $\geq$
  trained).  DNABERT-2 shows learned content attention but excludes
  ALiBi position biases ($^*$).}
\label{tab:attention}
\begin{tabular}{@{}lccl@{}}
\toprule
\textbf{Model} & \textbf{Trained} & \textbf{Untrained} &
  \textbf{Interpretation} \\
\midrule
NT~v2      & 0.322 & 0.328 & Architectural \\
DNABERT-2$^*$  & 0.257 & 0.000 & Learned (content-only) \\
RiNALMo    & 0.101 & 0.082 & Weak ($\Delta = +0.019$) \\
ERNIE-RNA  & 0.099 & 0.075 & Weak ($\Delta = +0.024$) \\
SpliceBERT & 0.064 & 0.056 & None ($\Delta = +0.008$) \\
RNA-FM     & 0.051 & 0.060 & None \\
UTR-LM     & 0.052 & 0.053 & None \\
\bottomrule
\end{tabular}
\end{table}

Among character-tokenized models, RiNALMo and ERNIE-RNA show the
strongest suggestions of learned attention: RiNALMo trained $\rho =
0.101$ versus untrained $0.082$ ($\Delta = +0.019$), ERNIE-RNA
trained $\rho = 0.099$ versus untrained $0.075$ ($\Delta = +0.024$).
Both deltas are small.  SpliceBERT (0.064 vs 0.056), RNA-FM (0.051
vs 0.060), and UTR-LM (0.052 vs 0.053) show no learned signal.
DNABERT-2's content-only attention (computed from Wqkv projections
without ALiBi position biases; Section~\ref{sec:phase4}) shows the
largest trained/untrained gap ($0.257$ vs $0.000$), but this metric
excludes the position-dependent biases that dominate the model's
actual attention computation.

NT~v2's 6-mer tokenization explains the architectural correlation.
Base-pairing partners---typically 3--30 nucleotides apart---are
aggregated into adjacent tokens, producing structured attention
patterns from positional proximity alone.  Random weights with this
tokenization already capture enough pairwise structure to match
trained attention.

\subsection{Scale does not compensate for domain mismatch}
\label{sec:scale}

Evo (7B parameters, byte-level tokenization, DNA-pretrained) produces
the weakest mutation sensitivity of all trained models
(Table~\ref{tab:phase1_summary}).  Two families (tRNA-Phe: 0.927,
tRNA-Ala: 0.899) show ratios below 1.0---Evo is less sensitive to stem
mutations than loop mutations for these structures.  Probing accuracy
(0.615) peaks at layer~0 (embedding) and falls below Caduceus (0.628,
14M) and HyenaDNA (0.621, 5.4M).  Despite being $500\times$ larger
than Caduceus and $70\times$ larger than RNA-FM, Evo shows no
advantage on any structure metric.  This observation is consistent
across mutation sensitivity, probing, and the null comparison, though
$n = 5$ models is too few to support a general scaling claim.

\begin{table}[t]
\centering
\caption{Mutation sensitivity summary, Phase~1 ($N = 12$).}
\label{tab:phase1_summary}
\begin{tabular}{@{}lccc@{}}
\toprule
\textbf{Model} & \textbf{Mean ratio} &
  \textbf{Families $>$ null} & \textbf{Probing $\Delta$} \\
\midrule
RNA-FM (99M, RNA)    & 1.774 & 1/12 & +0.013 \\
Caduceus (14M, DNA)  & 1.257 & 3/12 & +0.041 \\
HyenaDNA (5.4M, DNA) & 1.223 & 0/12 & +0.034 \\
Evo (7B, DNA)        & 1.170 & 0/12 & +0.028 \\
NT~v2 (56M, DNA)     & 1.094 & 6/12 & +0.016 \\
\bottomrule
\end{tabular}
\end{table}

\subsection{Composition controls absorb the embedding-level signal}
\label{sec:composition}

The nucleotide-stratified null reduced the apparent RNA-FM signal from
${\sim}5.5\times$ (naive, before composition control) to $1.8\times$
(composition-controlled).  The 24.7 percentage-point GC enrichment in
stems was the dominant contributor.  Per-family results
(Table~\ref{tab:per_family_phase1}) show that only one family
exceeded the first-order null for RNA-FM: tRNA-Ala (3.208), with
high-GC stems and low-GC loops where residual dinucleotide-context
effects contribute.  NT~v2 shows 6/12 families exceeding the null
despite having the lowest mean ratio (1.094), suggesting a more
consistent per-family signal than RNA-FM's high-mean, high-variance
pattern.

\begin{table}[t]
\centering
\caption{Per-family mutation sensitivity ratio, Phase~1 ($N = 12$).}
\label{tab:per_family_phase1}
\small
\begin{tabular}{@{}lcccccc@{}}
\toprule
\textbf{Family} & \textbf{RNA-FM} & \textbf{Untrained} &
  \textbf{NT~v2} & \textbf{HyenaDNA} & \textbf{Caduceus} &
  \textbf{Evo} \\
\midrule
tRNA-Phe (yeast)  & 2.758 & 1.050 & 1.180 & 1.185 & 1.374 & 0.927 \\
tRNA-Ala (human)  & 3.208 & 1.091 & 1.207 & 1.284 & 1.328 & 0.899 \\
5S rRNA           & 1.157 & 1.008 & 0.963 & 1.012 & 1.075 & 1.048 \\
Hammerhead        & 1.328 & 1.076 & 1.118 & 1.292 & 1.456 & 1.147 \\
SAM riboswitch    & 2.535 & 1.056 & 1.042 & 1.074 & 1.269 & 1.283 \\
TPP riboswitch    & 1.442 & 1.093 & 1.058 & 1.227 & 1.324 & 1.888 \\
SRP RNA           & 1.261 & 1.035 & 1.168 & 1.029 & 1.136 & 0.925 \\
miR-21            & 1.299 & 1.049 & 1.086 & 1.868 & 1.341 & 1.156 \\
IRES HCV dom.~II  & 1.710 & 1.045 & 0.928 & 1.190 & 1.225 & 0.994 \\
HDV ribozyme      & 1.135 & 1.021 & 1.175 & 1.150 & 1.291 & 1.316 \\
RNase~P           & 1.513 & 1.034 & 0.998 & 1.217 & 1.086 & 1.217 \\
U2 snRNA          & 2.200 & 1.024 & 1.203 & 1.144 & 1.173 & 1.234 \\
\midrule
\textbf{Mean}     & \textbf{1.774} & 1.048 & 1.094 & 1.223 & 1.257
                  & 1.170 \\
\textbf{$>$ null} & \textbf{1/12} & --- & 6/12 & 0/12 & 3/12 & 0/12 \\
\bottomrule
\end{tabular}
\end{table}

Phase~1 preregistered hypotheses confirmed this pattern.  H1 (trained/untrained
ratio $\geq 2.0$) failed at $1.69\times$.  H1b ($\geq 6/12$ families
exceed null) failed at 1/12.  H5 (probing $>$ baseline) failed, with
margins of +0.013 to +0.041 within noise at $N = 12$.  RNA-FM trained
did exceed untrained in 12/12 families (Wilcoxon $p = 0.000244$), but
this per-family unanimity was not preregistered and is reported as
exploratory.

\subsection{Phase~2 confirms and refines at $N = 52$}
\label{sec:phase2}

Phase~2 evaluates the original five models on 52 Rfam families under
the corrected max-over-layers null.  RNA-FM has the highest mean ratio
(1.774) but exceeds the null in only 1/52 families (2\%), barely above
the untrained control's 3/52 chance rate
(Table~\ref{tab:phase2_summary}).  The sole RNA-FM survivor is
tRNA-Ala ($3.208$ vs null $3.179$), a borderline exceedance absorbed
by the dinucleotide null (Section~\ref{sec:dinuc_results}).  All five
models fall \emph{below} the 58.1\% majority baseline on probing at
$N = 52$; the positive Phase~1 margins were noise amplified by small
sample size.

\begin{table}[t]
\centering
\caption{Phase~2 mutation sensitivity ($N = 52$).}
\label{tab:phase2_summary}
\begin{tabular}{@{}lccc@{}}
\toprule
\textbf{Model} & \textbf{Mean ratio} &
  \textbf{Families $>$ nuc null} & \textbf{Probing $\Delta$} \\
\midrule
RNA-FM (99M, RNA)         & 1.774 & 1/52  & $-0.010$ \\
NT~v2 (56M, DNA)          & 1.169 & 28/52 & $-0.039$ \\
Caduceus (14M, DNA)       & 1.197 & 5/52  & $-0.021$ \\
HyenaDNA (5.4M, DNA)      & 1.138 & 8/52  & $-0.058$ \\
Evo (7B, DNA)             & 1.352 & 6/52  & $-0.042$ \\
RNA-FM untrained          & 1.038 & 3/52  & --- \\
\bottomrule
\end{tabular}
\end{table}

H6 (ratio $\geq 2.0$) failed at $1.77\times$.  H6b ($\geq 9$ exceed
null) failed at 1/52.  H7 ($\geq 1$ family survives dinucleotide
null) failed: the sole first-order survivor (tRNA-Ala) is absorbed by
the dinucleotide null.  H8 (probing $>$ baseline $+ 0.02$) failed at
$-0.010$.  H10 (sign test) passed: RNA-FM trained exceeds untrained in
50 of 52 families (96\%), confirming the Phase~1 exploratory
observation as preregistered-confirmatory.  The two exceptions are
families with fewer than 10 stem positions where ratio estimates are
noisy.  The magnitude ($1.77\times$) remains below the 2.0 threshold
for practical claims.  H11 (NT~v2 attention) \textbf{failed}: trained $\rho =
0.322$ across 52 families, but untrained $= 0.328$ after correcting
for the SDPA backend bug that had produced a spurious $0.000$.  The
correlation is architectural (Section~\ref{sec:attn_results}).

\subsection{Dinucleotide null absorbs all survivors}
\label{sec:dinuc_results}

The dinucleotide-stratified null was run on the sole first-order
survivor (tRNA-Ala) plus SAM riboswitch (a Phase~1 survivor, included
regardless of Phase~2 result).  Both fall below the second-order
control (Table~\ref{tab:dinuc}).  tRNA-Ala ($3.208$) is absorbed by a
dinucleotide null of $4.070$, confirming that its signal reflected GC
dinucleotide enrichment rather than learned structure.  SAM riboswitch
($2.535$) is similarly absorbed ($3.600$).

\begin{table}[t]
\centering
\caption{Dinucleotide null results.  No family survives both
  composition controls.}
\label{tab:dinuc}
\begin{tabular}{@{}lcccc@{}}
\toprule
\textbf{Family} & \textbf{Ratio} &
  \textbf{Nuc 95th} & \textbf{Dinuc 95th} & \textbf{$>$ dinuc} \\
\midrule
tRNA-Ala          & 3.208 & 3.179 & 4.070 & No \\
SAM riboswitch    & 2.535 & 3.022 & 3.600 & No \\
\bottomrule
\end{tabular}
\end{table}

No RNA family survives both orders of composition control.  The
primary evidence for a learned structure effect is the directional
result (H10: 50/52 trained $>$ untrained), which does not depend on
any individual family exceeding any threshold.  The effect is
consistent but small ($1.77\times$), and entirely attributable to
sequence composition once stratified nulls are applied.

\subsection{Phase~3: RNA-specialized models}
\label{sec:phase3}

Phase~3 adds two RNA-pretrained models---RiNALMo (650M,
BERT+RoPE+SwiGLU) \citep{penic2024rinalmo} and UTR-LM (1.2M,
transformer encoder) \citep{chu2024utrlm}---bringing the total to
seven models across five architectures.  Both are evaluated on all 52
families with the same metrics as Phase~2, plus trained/untrained
comparisons for mutation sensitivity and attention.

\begin{table}[t]
\centering
\caption{Phase~3 results ($N = 52$).  RiNALMo shows a consistent
  directional signal (49/52 sign test) but a modest effect size.
  UTR-LM shows no structure sensitivity.}
\label{tab:phase3}
\begin{tabular}{@{}lcccccc@{}}
\toprule
\textbf{Model} & \textbf{Mean ratio} & \textbf{$>$ nuc null} &
  \textbf{Sign test} & \textbf{Probing} &
  \textbf{Attn $\rho$} & \textbf{Attn $\rho$ (unt.)} \\
\midrule
RiNALMo & 1.109 & 17/52 & 49/52 (94.2\%) & 0.561 & 0.101 & 0.082 \\
UTR-LM  & 1.045 & 3/52  & 37/52 (71.2\%) & 0.577 & 0.052 & 0.053 \\
\bottomrule
\end{tabular}
\end{table}

RiNALMo produces a higher mutation sensitivity ratio than any
DNA-pretrained model except RNA-FM (Table~\ref{tab:phase3}).  The
sign test (49/52 trained $>$ untrained, 94.2\%) exceeds the
preregistered H15 threshold of 75\%, confirming a directional effect.
However, the mean ratio (1.109) falls well below the 2.0 practical
threshold (H14: confirmed), and only 17/52 families exceed the
nucleotide-stratified null.  RiNALMo's attention-contact correlation
(0.101 trained vs 0.082 untrained) is the only trained-exceeds-untrained
attention result across all seven models, though the per-family sign
test (35/52, 67.3\%) falls short of strong directional evidence.

UTR-LM shows the weakest mutation sensitivity of all seven models
(1.045), with only 3/52 families exceeding the null---comparable to
the RNA-FM untrained baseline (3/52).  UTR-LM's 1.2M parameter count
and UTR-focused pretraining produce no detectable structure signal on
full-length RNA families.  Its probing accuracy (0.577) is the highest
of all models, consistent with the baseline probing result reflecting
embedding dimensionality rather than learned structure
(Section~\ref{sec:phase2}).  UTR-LM attention (0.052 trained vs 0.053
untrained) shows no learned structure.

H12 (RiNALMo attention $\rho < 0.10$) \textbf{rejected}: trained
$\rho = 0.101$, barely exceeding the threshold.  H13 (UTR-LM
attention $\rho > 0.15$) \textbf{rejected}: $\rho = 0.052$.  H14
(RiNALMo ratio $< 2.0$) \textbf{confirmed}: ratio $= 1.109$.  H15
(RiNALMo sign test $\geq 75\%$) \textbf{confirmed}: 49/52 (94.2\%).

\subsection{Phase~4: Expanded model coverage}
\label{sec:phase4}

Phase~4 adds ERNIE-RNA (86M, BERT, RNA-pretrained;
\citealp{wang2024ernierna}), SpliceBERT (19.4M, BERT,
RNA-pretrained; \citealp{chen2024splicebert}), and DNABERT-2 (117M,
MosaicBERT with ALiBi, DNA-pretrained; \citealp{zhou2024dnabert2}),
bringing the total to ten models.

\begin{table}[t]
\centering
\caption{Phase~4 results ($N = 52$).  ERNIE-RNA and SpliceBERT
  confirm the null pattern.  DNABERT-2 mutation and probing are not
  extractable due to BPE tokenization misalignment.}
\label{tab:phase4}
\begin{tabular}{@{}lcccccc@{}}
\toprule
\textbf{Model} & \textbf{Mean ratio} & \textbf{$>$ nuc null} &
  \textbf{Sign test} & \textbf{Probing} &
  \textbf{Attn $\rho$} & \textbf{Attn $\rho$ (unt.)} \\
\midrule
ERNIE-RNA  & 1.149 & 22/52 & 52/52 (100\%) & 0.604 & 0.099 & 0.075 \\
SpliceBERT & 1.045 & 4/52  & 48/52 (92.3\%) & 0.580 & 0.064 & 0.056 \\
DNABERT-2  & ---   & ---   & ---            & ---   & 0.257$^*$ & 0.000 \\
\bottomrule
\end{tabular}
\medskip

\small $^*$Content-only attention from Wqkv hooks; excludes ALiBi
position biases.
\end{table}

ERNIE-RNA produces the strongest sign test of all ten models: trained
exceeds untrained in 52/52 families (100\%), with a mean ratio of
1.149 (Table~\ref{tab:phase4}).  Despite this perfect directional
consistency, the mean ratio remains well below the 2.0 practical
threshold, and 22/52 families exceed the nucleotide null---a rate
driven by the consistent small effect rather than large individual
exceedances.  ERNIE-RNA's probing accuracy (0.604) is the highest of
all models, though this reflects embedding dimensionality rather than
structure discrimination.

SpliceBERT (19.4M) shows the weakest mutation sensitivity among
RNA-pretrained models (1.045), comparable to UTR-LM (1.045) and the
untrained baseline.  Only 4/52 families exceed the null.  The sign
test (48/52, 92.3\%) exceeds the 75\% threshold, confirming a
directional effect that is consistent but negligible in magnitude.

DNABERT-2 uses BPE tokenization and MosaicBERT architecture with
ALiBi position biases, which prevent standard metric extraction.  BPE
tokens do not align with nucleotide positions, making per-position
mutation sensitivity and probing undefined.  MosaicBERT's fused Wqkv
projection and ALiBi biases prevent standard attention extraction;
we compute content-only attention by hooking the Wqkv output,
splitting into Q and K, and computing $\mathrm{softmax}(QK^T / \sqrt{d_k})$
without ALiBi.  The resulting trained $\rho = 0.257$ versus untrained
$0.000$ shows learned content attention that correlates with base
pairing---unlike NT~v2, whose correlation is architectural.  The
untrained model produces uniform content attention (all Spearman
correlations zero) because random Wqkv weights yield unstructured
query-key products.  With ALiBi biases included, the true attention
pattern would be dominated by local position bias, so $0.257$ is an
upper bound on the full model's attention-contact correlation.

H16 (ERNIE-RNA attention $< 0.10$) \textbf{confirmed}: $\rho =
0.099$, a marginal pass.  H17 (ERNIE-RNA sign test $\geq 75\%$)
\textbf{confirmed}: 52/52 (100\%).  H18 (SpliceBERT attention $<
0.10$) \textbf{confirmed}: $\rho = 0.064$.  H19 (SpliceBERT ratio
$< 2.0$) \textbf{confirmed}: 1.045.  H20 (DNABERT-2
$|\text{trained} - \text{untrained}| < 0.05$) \textbf{rejected}:
$\Delta = 0.257$ (content-only).  H21 (DNABERT-2 trained $\rho <
0.20$) \textbf{rejected}: $\rho = 0.257$ (content-only).


%% ────────────────────────────────────────────
\section{Discussion}
\label{sec:discussion}

The composition confound we identified---GC enrichment in stems
inflating embedding-level structure metrics---is an instance of a
general problem: whenever the target structure correlates with feature
composition, any embedding metric sensitive to feature identity
inherits the correlation.  The nucleotide-stratified permutation null
is a template: stratify by the confounding feature, shuffle labels
within strata, and recompute the metric under the same selection
procedure used on real data.  The same principle applies to protein
secondary structure (amino acid composition differs between helices,
sheets, and coils), chromatin accessibility (GC content correlates
with open chromatin), and any biological structure correlated with
sequence composition.  The Johnson--Lindenstrauss lemma guarantees
that random projections of dimension $d \geq O(\epsilon^{-2} \log n)$
preserve pairwise distances with distortion $\leq \epsilon$
\citep{johnson1984extensions}, meaning any subspace-distance metric
computed on embeddings of dimension $d \gg \log n$ produces
non-trivial scores from random representations.  Composition controls
are a prerequisite for interpreting geometric structure metrics.

NT~v2's attention-contact correlation ($\rho = 0.322$) was initially
reported as fully learned (untrained $= 0.000$), a finding that
appeared to demonstrate structure encoding from DNA pretraining.  The
corrected analysis (untrained $= 0.328$) reverses this
interpretation: the correlation is an architectural property of the
ESM attention mechanism combined with 6-mer tokenization, present from
random initialization.  Among the seven character-tokenized attention
models (Table~\ref{tab:attention}), all show trained attention within
$\pm 0.025$ of untrained.  DNABERT-2 is the exception: its
content-only attention ($\rho = 0.257$ trained vs $0.000$ untrained)
shows learned structure-aware content patterns, though the practical
attention includes ALiBi position biases that dominate the
computation.  The SDPA bug that produced the spurious zero is a
cautionary example: HuggingFace's default attention backend silently
returns \texttt{None} when \texttt{output\_attentions=True}, and
downstream code that interprets absent attention as zero correlation
produces a qualitatively wrong result.

The two-order composition cascade validates the framework's design.
The first-order null absorbs ${\sim}3\times$ of the naive signal; the
second-order null absorbs what remains.  The Phase~1 survivors
tRNA-Ala and SAM riboswitch are both absorbed by second-order
controls.  No family survives both orders of composition control,
confirming that the embedding-level structure signal is entirely
attributable to nucleotide and dinucleotide composition.

The 2.0 threshold for H1/H6 was chosen to demand an effect large
enough for practical downstream use.  H6 fails while H10 passes,
separating magnitude from direction: RNA-FM's representations are
consistently more structure-sensitive than random weights across
50/52 families, but the $1.77\times$ effect size is insufficient for
applications requiring reliable stem-loop discrimination.

Ten models spanning attention-based, SSM-based, and hybrid
architectures provide broad coverage, though representation remains
thin per architecture.  Phase~4 extends the null attention result to
five RNA-pretrained BERT-style models: RNA-FM, RiNALMo, UTR-LM,
ERNIE-RNA, and SpliceBERT all show trained attention within $\pm
0.025$ of untrained.  ERNIE-RNA's mutation sign test (52/52, 100\%)
is the strongest directional result, followed by RNA-FM (50/52,
96.2\%) and RiNALMo (49/52, 94.2\%), consistent with RNA pretraining
conferring a small directional advantage that does not cross practical
thresholds.  SpliceBERT and UTR-LM show minimal mutation sensitivity
despite RNA pretraining, suggesting that model capacity (19.4M and
1.2M respectively) limits structure encoding.  DNABERT-2's BPE
tokenization prevents per-position metrics but illustrates a broader
limitation: subword tokenization breaks the nucleotide-level
assumptions of the evaluation framework.  The evaluation targets
secondary structure only and cannot assess tertiary contacts or
pseudoknots.  The observation that Evo underperforms smaller models
rests on a single model per size class.


%% ────────────────────────────────────────────
\section{Conclusion}
\label{sec:conclusion}

A composition-controlled evaluation framework reveals that GC
enrichment in stems inflates naive structure-sensitivity metrics by
${\sim}3\times$.  After controlling at one or two orders of sequence
composition, the residual embedding-level structure signal falls below
a preregistered threshold for practical significance.

No positive result survives all controls.  NT~v2's attention-contact
correlation ($\rho = 0.322$) is architectural: untrained weights
produce $\rho = 0.328$, and the earlier report of $0.000$ was an
artifact of the SDPA attention backend.  Seven character-tokenized
attention models show trained attention indistinguishable from
untrained; DNABERT-2 shows learned content attention ($\rho = 0.257$)
that excludes ALiBi position biases.  Phase~2 ($N = 52$) confirms:
the directional effect is real (50/52, H10 PASS), the magnitude is
modest (H6 FAIL at $1.77\times$), and composition controls absorb
every first-order survivor.  Phases~3--4 add five more models:
ERNIE-RNA achieves the strongest sign test (52/52, 100\%) with a mean
ratio of 1.149, RiNALMo confirms a directional signal (49/52) at
1.109, and SpliceBERT and UTR-LM show no structure sensitivity
($\leq 1.045$).

The framework---stratified permutation nulls, max-over-layers
matching, preregistered thresholds, untrained baselines---provides a
reusable template for evaluating structure awareness in sequence
foundation models.

\subsection*{Data availability}

Phase~1 results: \texttt{data/gpu\_results/batch3\_structure\_metrics/}.
Phase~2 results: Modal volume \texttt{prereg-experiment-results},
path \texttt{structure\_metrics/expanded\_rfam/}.
Phase~3 results: \texttt{data/gpu\_results/expanded\_rfam/}
(\texttt{rinalmo\_expanded\_*.json}, \texttt{utrlm\_expanded\_*.json},
\texttt{nt\_expanded\_20260713\_154855.json} with eager attention fix).
Phase~4 results: \texttt{data/gpu\_results/expanded\_rfam/}
(\texttt{ernierna\_expanded\_20260713\_193217.json},
\texttt{splicebert\_expanded\_20260713\_193419.json},
\texttt{dnabert2\_expanded\_20260713\_201316.json}).
RNA structures, analysis code, and preregistration documents are in
the \texttt{causal-rna} and \texttt{factorization-unified}
repositories.
Phase~1 prereg SHA: \texttt{694b43b}/\texttt{a7d10f5}.
Phase~2 prereg SHA: \texttt{bd4b3fd}/\texttt{74c8f49}.
Model checkpoints: RNA-FM (\texttt{ml4bio/RNA-FM}), NT~v2
(\texttt{InstaDeepAI/nucleotide-transformer-v2-50m-multi-species}),
HyenaDNA (\texttt{LongSafari/hyenadna-small-32k-seqlen}),
Caduceus (\texttt{kuleshov-group/caduceus-ps\_seqlen-131k\_%
d\_model-256\_n\_layer-16}),
Evo (\texttt{togethercomputer/evo-1-131k-base}),
RiNALMo (\texttt{ml4bio/RiNALMo}),
UTR-LM (\texttt{ml4bio/UTR-LM}),
ERNIE-RNA (\texttt{multimolecule/ernierna}),
SpliceBERT (\texttt{multimolecule/splicebert}),
DNABERT-2 (\texttt{zhihan1996/DNABERT-2-117M}).


%% ────────────────────────────────────────────
\bibliographystyle{plainnat}

\begin{thebibliography}{99}

\bibitem[Chen et al.(2022)]{rnafm}
Chen, J., Hu, Z., Sun, S., Tan, Q., Wang, Y., Yu, Q., Zong,
  L., Hong, L., Xiao, J., King, I., et~al.
\newblock Interpretable {RNA} foundation model from unannotated
  data for highly accurate {RNA} structure and function predictions.
\newblock \emph{arXiv:2204.00300}, 2022.

\bibitem[Dalla-Torre et al.(2023)]{dalla2023nucleotide}
Dalla-Torre, H., Gonzalez, L., Mendoza-Revilla, J.,
  Carranza, N.~L., Grzywaczewski, A.~H., et~al.
\newblock The {Nucleotide Transformer}: Building and evaluating
  robust foundation models for human genomics.
\newblock \emph{bioRxiv:2023.01.11.523679}, 2023.

\bibitem[Nguyen et al.(2023)]{nguyen2024hyenadna}
Nguyen, E., Poli, M., Faber, M., et~al.
\newblock {HyenaDNA}: Long-range genomic sequence modeling at
  single nucleotide resolution.
\newblock In \emph{NeurIPS}, 2023.

\bibitem[Schiff et al.(2024)]{schiff2024caduceus}
Schiff, Y., Kao, C.-H., Gokaslan, A., Dao, T., Gu, A.,
  Kuleshov, V.
\newblock Caduceus: Bi-directional equivariant long-range {DNA}
  sequence modeling.
\newblock In \emph{ICML}, 2024.

\bibitem[Nguyen et al.(2024)]{nguyen2024evo}
Nguyen, E., Poli, M., Durrant, M.~G., et~al.
\newblock Sequence modeling and design from molecular to genome
  scale with {Evo}.
\newblock \emph{Science}, 386(6723), 2024.

\bibitem[Johnson \& Lindenstrauss(1984)]{johnson1984extensions}
Johnson, W.~B. \& Lindenstrauss, J.
\newblock Extensions of {L}ipschitz mappings into a {H}ilbert space.
\newblock \emph{Contemporary Mathematics}, 26:189--206, 1984.

\bibitem[Tower(2026)]{tower2026boundary}
Tower, E.
\newblock Boundary conditions for geometric evaluation of
  foundation model representations.
\newblock \emph{Working paper}, 2026.

\bibitem[Ethayarajh(2019)]{ethayarajh2019}
Ethayarajh, K.
\newblock How contextual are contextualized word representations?
  {C}omparing the geometry of {BERT}, {ELMo}, and {GPT-2} embeddings.
\newblock In \emph{EMNLP}, 2019.

\bibitem[Vig et al.(2020)]{vig2020bertology}
Vig, J., Madani, A., Varshney, L.~R., Xiong, C., Socher, R.,
  \& Rajani, N.~F.
\newblock {BERTology} meets biology: Interpreting attention in
  protein language models.
\newblock In \emph{ICLR}, 2020.

\bibitem[Frazer et al.(2021)]{frazer2021disease}
Frazer, J., Notin, P., Dias, M., et~al.
\newblock Disease variant prediction with deep generative models
  of evolutionary data.
\newblock \emph{Nature}, 599(7883):91--95, 2021.

\bibitem[Karlin \& Bur\'{e}(1998)]{karlin1998dinucleotide}
Karlin, S. \& Bur\'{e}, C.
\newblock Dinucleotide relative abundance extremes: A genomic
  signature.
\newblock \emph{Trends in Genetics}, 14(10):398--402, 1998.

\bibitem[Peni\'{c} et al.(2024)]{penic2024rinalmo}
Peni\'{c}, R.~J., Vlah, D., Huber, R., Heider, D., \&
  \v{S}iki\'{c}, M.
\newblock {RiNALMo}: General-purpose {RNA} language models can
  generalize well on structure prediction tasks.
\newblock \emph{arXiv:2403.00043}, 2024.

\bibitem[Chu et al.(2024)]{chu2024utrlm}
Chu, Y., Yu, D., Li, Y., Huang, K., Shen, Y., Cong, L.,
  Zhang, J., \& Wang, M.
\newblock A 5' {UTR} language model for decoding untranslated
  regions of {mRNA} and function predictions.
\newblock \emph{Nature Machine Intelligence}, 6:449--460, 2024.

\bibitem[Wang et al.(2024)]{wang2024ernierna}
Wang, N., Bian, J., Li, J., et~al.
\newblock {ERNIE-RNA}: An {RNA} language model with structure-enhanced
  representations.
\newblock \emph{arXiv:2305.14196}, 2024.

\bibitem[Chen et al.(2024)]{chen2024splicebert}
Chen, K., Zhou, Y., Ding, M., Wang, Y., Ren, Z., \& Yang, Y.
\newblock Self-supervised learning on millions of pre-{mRNA} sequences
  improves sequence-based {RNA} splicing prediction.
\newblock \emph{Briefings in Bioinformatics}, 25(1), 2024.

\bibitem[Zhou et al.(2024)]{zhou2024dnabert2}
Zhou, Z., Ji, Y., Li, W., Dutta, P., Davuluri, R.~V., \& Liu, H.
\newblock {DNABERT-2}: Efficient foundation model and benchmark for
  multi-species genome.
\newblock In \emph{ICLR}, 2024.

\end{thebibliography}


%% ────────────────────────────────────────────
\appendix
\section{RNA family details ($N = 52$)}
\label{sec:appendix_families}

Phase~1 families (12, from original preregistration):

\begin{table}[h]
\centering
\small
\begin{tabular}{@{}lcccc@{}}
\toprule
\textbf{Family} & \textbf{Stem} & \textbf{Loop} &
  \textbf{RNA-FM ratio} & \textbf{$>$ nuc null} \\
\midrule
tRNA-Phe (yeast)    & 42 & 34 & 2.758 & No \\
tRNA-Ala (human)    & 42 & 34 & 3.208 & Yes \\
5S rRNA             & 72 & 48 & 1.158 & No \\
Hammerhead          & 16 & 32 & 1.328 & No \\
SAM riboswitch      & 52 & 42 & 2.535 & No \\
TPP riboswitch      & 40 & 40 & 1.442 & No \\
SRP RNA             & 44 & 33 & 1.261 & No \\
miR-21              & 64 &  8 & 1.299 & No \\
IRES HCV dom.~II    & 50 & 27 & 1.710 & No \\
HDV ribozyme        & 48 & 28 & 1.135 & No \\
RNase~P             & 42 & 36 & 1.513 & No \\
U2 snRNA            & 46 & 30 & 2.200 & No \\
\bottomrule
\end{tabular}
\end{table}

Expanded families (40 additional):

\begin{table}[h]
\centering
\small
\begin{tabular}{@{}llcccc@{}}
\toprule
\textbf{Family} & \textbf{Class} & \textbf{Stem} & \textbf{Loop} &
  \textbf{RNA-FM ratio} & \textbf{$>$ nuc null} \\
\midrule
6S RNA              & ncRNA       & 42 & 34 & 1.823 & No \\
7SK RNA             & ncRNA       & 38 & 31 & 1.593 & No \\
Bacterial SRP       & ncRNA       & 43 & 22 & 2.179 & No \\
c-di-GMP ribosw.    & Riboswitch  & 36 & 27 & 1.826 & No \\
Cobalamin ribosw.   & Riboswitch  & 38 & 36 & 2.214 & No \\
Corona 5' UTR       & Cis-reg     & 51 & 33 & 1.252 & No \\
Corona s2m          & Cis-reg     & 56 & 27 & 1.268 & No \\
CRISPR leader       & CRISPR      & 25 & 30 & 1.121 & No \\
CrPV IRES           & IRES        & 40 & 24 & 1.575 & No \\
Fluoride ribosw.    & Riboswitch  & 36 & 28 & 1.473 & No \\
FMN riboswitch      & Riboswitch  & 52 & 26 & 1.889 & No \\
glmS ribozyme       & Ribozyme    & 44 & 38 & 1.252 & No \\
Glycine ribosw.     & Riboswitch  & 26 & 32 & 1.167 & No \\
Group~I intron P4P6 & Intron      & 46 & 31 & 1.827 & No \\
Group~II intron D5  & Intron      & 38 & 21 & 1.011 & No \\
Hatchet ribozyme    & Ribozyme    & 26 & 31 & 1.865 & No \\
HCV IRES III        & IRES        & 40 & 24 & 1.480 &
  No \\
Histone 3' stem     & Cis-reg     & 10 & 10 & 4.078 & No \\
IRE stem-loop       & Cis-reg     & 21 &  9 & 2.006 & No \\
Lysine ribosw.      & Riboswitch  & 40 & 38 & 1.581 & No \\
Manganese ribosw.   & Riboswitch  & 40 & 24 & 1.008 & No \\
miR-122             & miRNA       & 44 & 28 & 1.252 & No \\
miR-155             & miRNA       & 44 & 22 & 2.529 & No \\
miR-let-7           & miRNA       & 48 & 27 & 0.952 & No \\
Pistol ribozyme     & Ribozyme    & 28 & 38 & 1.673 & No \\
preQ1 ribosw.       & Riboswitch  & 26 & 25 & 2.757 & No \\
Purine ribosw.      & Riboswitch  & 36 & 21 & 2.353 & No \\
SAH riboswitch      & Riboswitch  & 32 & 28 & 1.252 & No \\
SECIS element       & Cis-reg     & 30 & 28 & 1.166 & No \\
T-box leader        & Cis-reg     & 38 & 26 & 2.651 & No \\
THF riboswitch      & Riboswitch  & 52 & 30 & 1.150 & No \\
tmRNA               & ncRNA       & 56 & 29 & 2.682 & No \\
Twister ribozyme    & Ribozyme    & 30 & 37 & 1.584 & No \\
U1 snRNA            & snRNA       & 44 & 31 & 1.721 & No \\
U4 snRNA            & snRNA       & 38 & 34 & 2.101 & No \\
U5 snRNA            & snRNA       & 44 & 18 & 2.864 & No \\
U6 snRNA            & snRNA       & 40 & 30 & 1.091 & No \\
Vault RNA           & ncRNA       & 44 & 29 & 1.264 & No \\
Y RNA               & ncRNA       & 36 & 32 & 1.979 & No \\
ZMP riboswitch      & Riboswitch  & 30 & 26 & 1.961 & No \\
\bottomrule
\end{tabular}
\end{table}

\TODO{Add Rfam accession IDs and GC\% columns from family metadata.}

\end{document}
